\section{Mixed norms and the exact relative graph bound}\label{inf:sec-mixed}

The intermediate space keeps the radial variable Hilbertian and the
angular coefficients summable and preserves radial spectral
projections in endpoint-normalized angular coordinates.

\begin{definition}[Mixed space and weighted matrix norms]
\label{inf:mixed-spaces}
Let $\Mmix$ be the space of all families $(u_j)_{j\ge0}$ with
$u_j\in L^2((0,1),pt^{\beta_j}dt)$ and finite norm below,
represented as $u=\sum_jS_ju_j(t)$. Put
\begin{equation}\label{inf:mixed-norm}
 \nrm u_{\Mmix}=\sum_j2^j\nrm{u_j}_{\beta_j},\qquad
 \nrm v_\beta^2=p\int_0^1t^\beta|v(t)|^2\,dt.
\end{equation}
The actual invariant Hilbert space $\Hc$ has norm
\begin{equation}\label{inf:true-H-norm}
 \nrm u_\Hc^2=\sum_j\eta_j^2\nrm{u_j}_{\beta_j}^2.
\end{equation}
For an angular block matrix $T$ define, at fixed $\sigma>0$,
\begin{align}
 \nrm T_{\Mmix,\sigma}
 &=\sup_j\sum_i2^{i-j}e^{\sigma|i-j|}
                            \nrm{T_{ij}}_{\beta_j\to\beta_i},
 \label{inf:M-matrix-norm}\\
 \nrm T_{\Hc,\sigma}
 &=\max\left\{
 \sup_i\sum_je^{\sigma|i-j|}\nrm{T^H_{ij}},
 \sup_j\sum_ie^{\sigma|i-j|}\nrm{T^H_{ij}}\right\},
 \quad T^H_{ij}=\frac{\eta_i}{\eta_j}T_{ij}.
 \label{inf:H-matrix-norm}
\end{align}
The block norms in the second line use the same radial Hilbert spaces;
only the angular basis has been normalized.
\end{definition}

\begin{lemma}[Embedding and matrix algebra]\label{inf:mixed-algebra}
There are norm-at-most-one embeddings
$\Cc\hookrightarrow\Mmix\hookrightarrow\Hc$.
The space $\Mmix$ is complete, and finite polynomial sums are
dense in all three spaces.
The two matrix norms in
\eqref{inf:M-matrix-norm}--\eqref{inf:H-matrix-norm} are
submultiplicative, have identity norm one, and dominate their
respective operator norms. The spaces of matrices with finite such
norms are complete. For a block-diagonal operator both reduce
to the supremum of its radial block norms.
\end{lemma}
\begin{proof}
The radial squared basis norm is
$p/(p+2j+2s)\le1$. The triangle inequality gives the first
embedding from \eqref{inf:C-norm}; the inequality $\eta_j\le1$
and $\ell^2\le\ell^1$ give the second.
Finite angular truncations and polynomial density for each positive
radial weight give density in $\Mmix$ and $\Hc$; density in
$\Cc$ is part of its definition. Completeness of $\Mmix$ follows
from its description as a weighted $\ell^1$ sum of the complete
radial Hilbert spaces: a Cauchy family has a limit in every
radial block, and Fatou's inequality bounds the sum of the
blockwise differences by its Cauchy bound.
For products of blocks use
$e^{\sigma|i-j|}\le e^{\sigma|i-k|}e^{\sigma|k-j|}$ and
$2^{i-j}=2^{i-k}2^{k-j}$, then sum the positive block norm bounds.
The ratios $\eta_i/\eta_j$ telescope through the intermediate
index. This proves submultiplicativity for both row and column
sums. The mixed operator bound follows from summing columns;
the Hilbert bound is the block Schur inequality, obtained by
Cauchy--Schwarz with these positive row and column sums.
For completeness, a Cauchy sequence in either matrix norm is
Cauchy in every block operator norm. Take its blockwise limits.
Fatou's inequality for the positive row and column sums bounds
the norm of each difference from that limit by the corresponding
Cauchy bound, uniformly over the exterior index. The sequence
therefore converges in the matrix norm, and the preceding operator
bounds identify the limit with a bounded operator on the indicated
space. For diagonal blocks all off-diagonal weights are one and
only one term remains in each sum.
\end{proof}

\begin{lemma}[Solid multipliers in the two matrix norms]
\label{inf:mixed-multipliers}
For integers $i,h\ge0$,
\begin{equation}\label{inf:mixed-multiplier-bounds}
 \nrm{M_{S_it^h}}_{\Mmix,\sigma}\le2^ie^{\sigma i},\qquad
 \nrm{M_{S_it^h}}_{\Hc,\sigma}\le(2i+1)e^{\sigma i}.
\end{equation}
\end{lemma}
\begin{proof}
An angular product output has factor
$t^{i+j-k+h}u_j$. Its squared radial weight has exponent
\[
 \beta_k+2(i+j-k+h)=\beta_j+2i+2h\ge\beta_j.
\]
Multiplication by that power of $t$ is a contraction between the
indicated radial spaces. The positive angular coefficients sum to
one, and $|k-j|\le i$; this proves the mixed bound.
In the true Hilbert space multiplication has norm at most one
because $|S_it^h|\le1$ on the ball. Every angular block is an
orthogonal compression and has norm at most one. Its bandwidth
is $i$, so each weighted row and column contains at most $2i+1$
terms, proving the Hilbert bound.
\end{proof}

\begin{lemma}[Poisson and gradient estimates in mixed and Hilbert norms]
\label{inf:mixed-Poisson}
On all source layers one has
\begin{equation}\label{inf:mixed-Poisson-bounds}
 \nrm{\KD}\le Cp^{-2},\qquad
 \nrm{\Euler\KD}\le Cp^{-1},\qquad
 \nrm{\Euler^2\KD}\le C
\end{equation}
in both matrix norms. For $A=2i>0$,
\begin{align}
 \nrm{\nabla S_i\cdot\nabla\KD}_{\Mmix,\sigma}
 &\le CA2^ie^{\sigma i}/p,&
 \nrm{\nabla S_i\cdot\nabla\Euler\KD}_{\Mmix,\sigma}
 &\le CA2^ie^{\sigma i},\label{inf:mixed-gradients}\\
 \nrm{\nabla S_i\cdot\nabla\KD}_{\Hc,\sigma}
 &\le CA(2i+1)e^{\sigma i}/p,&
 \nrm{\nabla S_i\cdot\nabla\Euler\KD}_{\Hc,\sigma}
 &\le CA(2i+1)e^{\sigma i}.
 \label{inf:Hilbert-gradients}
\end{align}
The constants are absolute for fixed $\sigma$.
\end{lemma}
\begin{proof}
In angular degree $B=2j$, let $\kf$ solve
$4(t\kf''+(\beta+1)\kf')=f$, $\kf(1)=0$, $\beta=B+p-1$.
The radial spectral lower bound gives
$\nrm\kf_\beta\le C(B+p)^{-2}\nrm f_\beta$.
Multiplication of the equation by $\overline\kf t^\beta$ and integration
by parts give
\[
 4\nrm{\kf'}_{\beta+1}^2
 =-\operatorname{Re}\langle f,\kf\rangle_\beta,
 \qquad
 \nrm{\kf'}_{\beta+1}\le C(B+p)^{-1}\nrm f_\beta.
\]
These equalities first hold for polynomial sources; their
extensions follow from the Dirichlet form and Poisson uniqueness.
The profile identities
$\Euler\KD=B\kf+2t\kf'$ and
$\Euler^2\KD=B^2\kf+4(1-p)t\kf'+tf$
then give \eqref{inf:mixed-Poisson-bounds} on each block, hence in
both matrix norms.

For the gradient operators use the two profiles
\eqref{inf:gradient-profile-one}--\eqref{inf:gradient-profile-two}
with the general source $f$ in place of $J_s^\beta$ and
$\kf$ in place of $k_s$.
When $h\ge1$, the output squared weight for $t^{h-1}\kf$ is
$t^{\beta+A-2}$, which is bounded above by $t^\beta$ since
$A\ge2$. The corresponding weight for $t^h\kf'$ is
$t^{\beta+A}\le t^{\beta+1}$.
The source term $t^hf$ is also a contraction between its two
radial spaces. Equation~\eqref{inf:hZ-bound} and the preceding
energy bounds give $CA/(B+p)$ for the first profile and $CA$ for
the second. There is no negative-power term when $h=0$.
Positive angular summation and the bandwidth weight prove
\eqref{inf:mixed-gradients}.

In Hilbert coordinates, homogeneity and
Lemma~\ref{inf:angular-sup} give
$\nrm{\nabla S_i}_\infty\le\sqrt2A$.
For $u=\KD f$, the Dirichlet spectral lower bound gives
$\nrm{\nabla u}_2\le(p-1)^{-1}\nrm f_2$.
On the ball one also has
\begin{equation}\label{inf:convex-Hessian}
 \nrm{D^2u}_2\le\nrm{\Delta u}_2=\nrm f_2.
\end{equation}
For a smooth real Dirichlet function, integrate the divergence of
$(\Delta u)\nabla u-D^2u\nabla u$.
Its interior divergence is $(\Delta u)^2-|D^2u|^2$.
At the boundary $\nabla u=(\partial_\nu u)\nu$ and
$\Delta u=\partial_{\nu\nu}u+H_{\partial B}\partial_\nu u$,
so the boundary integral is
$\int_{\partial B}H_{\partial B}|\partial_\nu u|^2\ge0$.
This proves \eqref{inf:convex-Hessian}; use real and imaginary
parts and $H^2$ approximation for the general case.
Since $\nabla\Euler u=\nabla u+(D^2u)y$, the unweighted Hilbert
bounds for the two gradient operators are $CA/p$ and $CA$.
Their angular bandwidth is $i$, as is also visible in the exact
profiles. Compressing to blocks and summing at most $2i+1$
weighted entries per row and column proves
\eqref{inf:Hilbert-gradients}. This argument covers the harmonic
source layer as well as all the others.
\end{proof}

\begin{lemma}[Closed mixed graph realization]\label{inf:mixed-domain}
Let $H_0$ have its Hilbert Dirichlet realization. Its realization
on $\Mmix$ has domain
\[
 \Dom_{\Mmix}(H_0)=\KD\Mmix
 =\{u\in\Dom_{\Hc}(H_0):H_0u\in\Mmix\}.
\]
It is closed in $\Mmix$ and its graph norm is equivalent, for
fixed $p$, to the norm of its source. Its realization on $\Cc$
has the identical description with $\Cc$ in place of $\Mmix$.
Let $V$ be defined on $\Dom_{\Hc}(H_0)$, and suppose
$\nrm{VH_0^{-1}}<1$ both on $\Hc$ and on $\Mmix$.
Then $H_0+V$ restricts to a closed operator on
$\Dom_{\Mmix}(H_0)$, and
\[
 \{u\in\Dom_{\Hc}(H_0):(H_0+V)u\in\Mmix\}
 =\Dom_{\Mmix}(H_0).
\]
The same conclusions hold with $\Cc$ in place of $\Mmix$ when
the relative norm is below one on both $\Hc$ and $\Cc$.
\end{lemma}
\begin{proof}
The Poisson estimates give a bounded map $\KD:\Mmix\to\Mmix$,
and the Hilbert identification gives $H_0\KD f=-f$.
Thus every element in its range is in the stated Hilbert domain
and has source in $\Mmix$.
Conversely, if $u$ lies in that Hilbert domain and $H_0u=f\in\Mmix$,
Dirichlet uniqueness gives $u=-\KD f$, proving equality of the
descriptions. The norm $\nrm{H_0u}_{\Mmix}$ is a complete norm
on the domain, since $\KD$ is a bijection from the complete
source space onto it. The Poisson bound makes it equivalent to
$\nrm u_{\Mmix}+\nrm{H_0u}_{\Mmix}$.
If $u_n\to u$ and $H_0u_n\to f$ in $\Mmix$, the embeddings
imply the same convergence in $\Hc$. Closedness of the Hilbert
operator gives $u\in\Dom_\Hc(H_0)$ and $H_0u=f$; hence the
description just proved puts $u$ in the mixed domain. This
proves closedness directly.
The coefficient-space proof uses the coefficient Poisson bound
and the same Hilbert embedding, so every step remains valid.

In each indicated source space, the factor
$I+VH_0^{-1}$ is bounded and boundedly invertible in the source
space. The product
$(H_0+V)=(I+VH_0^{-1})H_0$ therefore has an equivalent graph norm
on the same complete domain and is closed. Adding a bounded
scalar shift preserves closedness there.
For a source in $\Mmix$, the partial sums of the Neumann series
for $(I+VH_0^{-1})^{-1}$ agree with the Hilbert partial sums.
Convergence in $\Mmix$ implies convergence in $\Hc$, so the
two inverses coincide on these sources. If
$u\in\Dom_{\Hc}(H_0)$ and $(H_0+V)u=f\in\Mmix$, Hilbert
uniqueness gives
$u=H_0^{-1}(I+VH_0^{-1})^{-1}f\in\Dom_{\Mmix}(H_0)$.
The reverse inclusion follows by applying the mixed-space
operator. This proves the displayed equality and the restriction
property. Replacing $\Mmix$ by $\Cc$ gives the coefficient
version. The same argument applies to
a bounded angular-diagonal radial spectral projection, which
commutes with $H_0$ and preserves the graph domain.
\end{proof}

\subsection{A centre-size bound independent of the accuracy order}

\begin{lemma}[Order-independent leading size]\label{inf:uniform-centre-size}
For a fixed centre order $N$, expand $f_N=\sum_j f_{N,j}G_j$ and put
\begin{equation}\label{inf:strong-centre-norm}
 \delta=\sum_j(1+j)^{10}(2e^2)^j\wt(2j)|f_{N,j}|.
\end{equation}
There is an absolute $C_0$ such that for every fixed $N$,
\begin{equation}\label{inf:uniform-delta}
 \delta\le C_0\tau p^{-2}\qquad(p\ge P_N).
\end{equation}
Only the threshold depends on $N$. This stronger norm is used for
the known finite centre; the unknown Newton shape retains the
space $\Ec_1$ with one moment and angular radius two.
\end{lemma}
\begin{proof}
Equation~\eqref{inf:centre-profile} gives
\[
 f_N=-\frac{\theta}{2R^2pg_2}\mathsf H_2+O_N(\tau p^{-4})
\]
in every fixed polynomial coefficient norm.
The leading term has only index two. Its coefficient is bounded
by $C\tau p^{-2}$ independently of $N$, since $pg_2>9/16$ and
$\mathsf H_2(1)\le1$. Its weight in \eqref{inf:strong-centre-norm} is an
absolute constant for $p\ge64$.
For the fixed jet the remainder has bounded degree depending on
$N$, and the finite monic-to-$G_j$ connection is uniformly
bounded for large $p$, because
$\mathsf H_j(1)\to(1-r)^j>0$ for each fixed $j$.
Thus its strong norm is at most $C_N\tau p^{-4}$.
Choose $C_0$ to be the leading bound plus one, and increase $P_N$
until $C_Np^{-2}\le1$. This proves \eqref{inf:uniform-delta}.
\end{proof}

\subsection{Scalar perturbation and density conjugation}

Use the fixed centre chart $\Phi_N=qy$, $d=q+\Euler q$, and define
\begin{equation}\label{inf:density-operators}
 H_\Phi=-P_{\Phi_N}\Delta,\qquad
 m=(\det D\Phi_N)^{1/2}=q^{p-1/2}d^{1/2},\qquad
 \Hu=mH_\Phi m^{-1},\qquad \mathcal V=\Hu-H_0.
\end{equation}
Here $\Hu$ is the unitary Dirichlet realization on $\Hc$.
The target volume uses the same fixed normalization as unit-ball
volume, so multiplication by $m$ is exactly the unitary density
factor. The scalar graph operator is
$J_D=(P_{\Phi_N}\Delta+R^2)\KD$.
All full-source graph estimates in this section are at this
fixed finite centre with the strong norm
\eqref{inf:strong-centre-norm}. The Newton iteration uses its
frozen inverse, while the nonlinear map is estimated on
$\Gc_p\times\Ec_1$.

\begin{proposition}[The Hilbert relative estimate]
\label{inf:Hilbert-relative}
For the real finite-centre chart in \eqref{inf:density-operators},
let $\delta$ be \eqref{inf:strong-centre-norm}. There are absolute
constants $C,\varepsilon_H>0$ such that, if $p\ge64$ and
$p\delta\le\varepsilon_H$, then
\begin{equation}\label{inf:Hilbert-relative-bound}
 \nrm{\mathcal VK_0^{-1}}_{\Hc,1}
 \le C(\delta+p\delta^2).
\end{equation}
The scalar graph difference also satisfies
\begin{equation}\label{inf:Hilbert-scalar-bound}
 \nrm{(P_{\Phi_N}\Delta-\Delta)\KD}_{\Hc,1}
 \le C(\delta+\delta^2).
\end{equation}
These constants do not depend on $N$, the split, or the number
of nonzero centre modes. All invariant radial source layers
are included.
\end{proposition}
\begin{proof}
Use the orthonormal angular coordinates of \eqref{inf:H-matrix-norm}.
Its ingredients are the endpoint supremum bound, finite angular
bandwidth, and ball Hilbert estimates.

\emph{1. Banded operators.}
Let $P_j^H$ be the orthogonal projection onto angular degree
$2j$. If an operator $T$ has angular bandwidth $D$, each block
$P_i^HTP_j^H$ has norm at most $\nrm T_{\Hc\to\Hc}$, and each
row and column has at most $2D+1$ nonzero blocks. Therefore
\begin{equation}\label{inf:Hilbert-band-bound}
 \nrm T_{\Hc,1}\le(2D+1)e^D\nrm T_{\Hc\to\Hc}.
\end{equation}
This bounds both weighted row and column sums.

\emph{2. Multipliers.}
For an invariant polynomial multiplier $a$ of Jacobi degree
(maximum angular index) at most $D$, orthogonality gives bandwidth $D$ and multiplication
on $L^2$ has norm at most $\nrm a_\infty$. Hence
\[
 \nrm{M_a}_{\Hc,1}\le(2D+1)e^D\nrm a_\infty.
\]
In particular $|S_it^h|\le1$ gives
$(2i+1)e^i$ for that multiplier. Summing the centre coefficients
shows that multiplication by $\mathfrak h=q-1$, $\Euler\mathfrak h$,
and $\Euler(\Euler-1)\mathfrak h$ has norm at most $C\delta$.
The fixed moments in \eqref{inf:strong-centre-norm} absorb all
the displayed degree factors. The endpoint bound used here is
for $\ab\le\bb$ and $G_i(1)=1$, as fixed in
\eqref{inf:angular-jacobi}.

\emph{3. Products of gradients.}
For $A=2i$ the solid harmonic has
$|\nabla S_i|\le\sqrt2A$, by radial homogeneity and the
Bernstein bound on great circles. Thus
$|\nabla S_i\cdot\nabla S_j|\le2(2i)(2j)$.
Its angular bandwidth as a multiplier is at most $i+j$.
Equation~\eqref{inf:Hilbert-band-bound} gives
\begin{equation}\label{inf:Hilbert-gradient-product}
 \nrm{M_{\nabla S_i\cdot\nabla S_j}}_{\Hc,1}
 \le2(2i)(2j)(2i+2j+1)e^{i+j}.
\end{equation}
After summing the shape coefficients, this bounds the multipliers
$M_q=\nabla q\cdot\nabla q$,
$\nabla q\cdot\nabla\Euler q$,
$\nabla d\cdot\nabla d$, and
$\nabla q\cdot\nabla d$ by $C\delta^2$.
An Euler derivative on a shape adds only its angular degree,
which is included in the ten moments. This Hilbert multiplier
bound is $C\delta^2$.

\emph{4. Diagonal Poisson operators.}
For $u=\KD f$, the ball spectral estimate and its energy identity
give
\[
 \nrm u_2\le(p-1)^{-2}\nrm f_2,\qquad
 \nrm{\nabla u}_2\le(p-1)^{-1}\nrm f_2.
\]
The boundary identity \eqref{inf:convex-Hessian} gives
$\nrm{D^2u}_2\le\nrm f_2$ on the convex ball.
Since $\Euler u=y\cdot\nabla u$ and
$\Euler^2u=y^TD^2uy+y\cdot\nabla u$, the three Poisson operators
have norms at most $(p-1)^{-2}$, $2/(p-1)$ and $4$.
They are angular diagonal, so the same bounds hold in
$\nrm\cdot_{\Hc,1}$ on all source layers.

\emph{5. Gradient--Poisson operators.}
The preceding energy estimate gives
$\nrm{\nabla S_i\cdot\nabla\KD}_{\Hc\to\Hc}
\le\sqrt2(2i)/(p-1)$.
Using $\nabla\Euler u=\nabla u+(D^2u)y$ gives
$\nrm{\nabla S_i\cdot\nabla\Euler\KD}_{\Hc\to\Hc}
\le\sqrt2(2i)(1+(p-1)^{-1})$.
Their exact profiles \eqref{inf:gradient-profile-one}--
\eqref{inf:gradient-profile-two} show angular bandwidth $i$.
Applying \eqref{inf:Hilbert-band-bound} and summing the centre
coefficients proves
\[
 \nrm{\nabla\mathfrak h\cdot\nabla\KD}_{\Hc,1}\le C\delta/p,
 \qquad
 \nrm{\nabla\mathfrak h\cdot\nabla\Euler\KD}_{\Hc,1}\le C\delta.
\]
The same estimates hold with $\mathfrak h$ replaced by
$\Euler\mathfrak h$, and hence for
$d-1=\mathfrak h+\Euler\mathfrak h$.

\emph{6. Inverse multipliers and density.}
The Hilbert matrix norm is a Banach algebra with unit norm one
by Lemma~\ref{inf:mixed-algebra}.
For $C\delta<1$, the Neumann series for $q^{-1},d^{-1}$
converge in it; they also converge uniformly on the ball to the
pointwise multipliers. Their fixed inverse powers are bounded.
The absolute binomial majorant for $m^{\pm1}$ gives
\begin{equation}\label{inf:Hilbert-density-size}
 \nrm{M_{m^{\pm1}}}_{\Hc,1}
 \le(1-C\delta)^{-p-1}\le2
\end{equation}
when the absolute constant $\varepsilon_H$ is sufficiently small.
This choice is independent of the finite centre order.

\emph{7. The scalar graph difference.}
Compose \eqref{inf:scalar-pullback} with $\KD$ and subtract
the identity. The first term costs $C\delta$ by the inverse
multiplier bound, and the second costs $C\delta$ by step 5.
The third costs $C\delta^2$ by steps 3 and 4.
Formula~\eqref{inf:Cq-identity} bounds its contracted Hessian
multiplier by $C\delta^2(1+\delta)$; it is followed by
$\Euler\KD$, of norm $C/p$. The remaining gradient-square
term is bounded in the same way. This proves
\eqref{inf:Hilbert-scalar-bound}.

\emph{8. The density commutator.}
Both $q$ and $d$ are harmonic. With $\alpha=p-1/2$,
\[
 \frac{\nabla m^{-1}}{m^{-1}}
 =-\alpha q^{-1}\nabla q-\tfrac12d^{-1}\nabla d,
\]
\[
 \frac{\Delta m^{-1}}{m^{-1}}
 =\alpha(\alpha+1)q^{-2}\nabla q\cdot\nabla q
 +\tfrac34d^{-2}\nabla d\cdot\nabla d
 +\alpha q^{-1}d^{-1}\nabla q\cdot\nabla d.
\]
The gradient part of $[\Delta,m^{-1}]\KD$ is bounded by
$Cp(\delta/p)=C\delta$, using step 5 and the inverse multipliers.
Its potential part is bounded by
$Cp^2\delta^2(p-1)^{-2}\le C\delta^2$, using steps 3 and 4.
Thus
\begin{equation}\label{inf:Hilbert-commutator}
 \nrm{[\Delta,m^{-1}]\KD}_{\Hc,1}
 \le C(\delta+\delta^2).
\end{equation}
The potential starts at quadratic order because $q,d$ are harmonic.

\emph{9. Assembly on the graph domain.}
Put $X=(H_\Phi-H_0)H_0^{-1}
=(P_{\Phi_N}\Delta-\Delta)\KD$.
Multiplication by $m^{-1}$ preserves the Dirichlet boundary
condition, and the preceding commutator estimate bounds
$H_0m^{-1}H_0^{-1}$ by an absolute constant. Hence
\[
 \mathcal VH_0^{-1}
 =mX(H_0m^{-1}H_0^{-1})+m[H_0,m^{-1}]H_0^{-1}
\]
has Hilbert weighted norm at most $C(\delta+\delta^2)$,
which is at most $C(\delta+p\delta^2)$.
Finally $H_0K_0^{-1}$ is an angular-diagonal radial spectral
contraction, so multiplying by it proves
\eqref{inf:Hilbert-relative-bound}.
The identities hold first on the smooth Dirichlet graph core;
the bounds extend them to the Hilbert graph realization.
\end{proof}

\begin{theorem}[Exact relative graph bounds]\label{inf:relative-graph}
Fix $\sigma=1$. There is an absolute $C$ independent of $N$ such
that, for every fixed centre order $N$ and all $p\ge P_N$,
\begin{align}
 \nrm{\mathcal VK_0^{-1}}_{\Mmix,1}
       +\nrm{\mathcal VK_0^{-1}}_{\Hc,1}
 &\le C\tau p^{-2},\qquad K_0=H_0+R^2,
 \label{inf:relative-two-norms}\\
 \nrm{(P_{\Phi_N}\Delta-\Delta)\KD}_{\Cc\to\Cc}
       +\nrm{(P_{\Phi_N}\Delta-\Delta)\KD}_{\Mmix\to\Mmix}
 &\le C\tau p^{-2}.\label{inf:relative-scalar}
\end{align}
The density multipliers and their graph conjugates also satisfy,
for the same $p\ge P_N$,
\begin{align}
 \max\{\nrm{M_{m^{\pm1}}}_{\Mmix,1},
                  \nrm{M_{m^{\pm1}}}_{\Hc,1}\}&\le2,\nonumber\\
 \max\{\nrm{H_0M_{m^{-1}}H_0^{-1}}_{\Mmix,1},
                  \nrm{H_0M_{m^{-1}}H_0^{-1}}_{\Hc,1}\}&\le C.
 \label{inf:density-graph-bounds}
\end{align}
All these bounds include every source layer. The operator $\Hu$ has domain
$H^2(B)\cap H^1_0(B)$ in its invariant sector and has the actual
centre-domain spectrum.
\end{theorem}
\begin{proof}
Proposition~\ref{inf:Hilbert-relative} supplies the Hilbert half.
For the mixed and coefficient bounds, use endpoint-normalized
angular coordinates.
We keep the size $\delta$ in \eqref{inf:strong-centre-norm}
throughout the calculation. In $\Cc$ and both weighted matrix
algebras, solid multipliers and their first two Euler derivatives
are controlled by their coefficient sums with the corresponding
degree factors. The ten moments and the factor $(2e^2)^j$
dominate the angular weights, the bandwidth factor $2j+1$, and
these fixed derivatives. Proposition~\ref{inf:module} and
Lemmas~\ref{inf:mixed-multipliers} and \ref{inf:mixed-Poisson} thus give
\[
 \nrm{M_{q-1}}+\nrm{M_{d-1}}\le C\delta,
 \quad
 \nrm{\nabla(q-1)\cdot\nabla\KD}\le C\delta/p,
 \quad
 \nrm{\nabla(q-1)\cdot\nabla\Euler\KD}\le C\delta.
\]
For the coefficient-space version use
Lemma~\ref{inf:gradient-profiles}, which explicitly includes $s=0$.
The same statements hold with a fixed Euler derivative on the
harmonic shape.

Products of two harmonic gradients have multiplier norm
$Cp\delta^2$. In the mixed and coefficient spaces this follows
from the exact product coefficient $2h(h+2k+p-1)c_{ij}^k$
and its bound by $Cp(1+i)(1+j)$.
In the Hilbert matrix algebra use the pointwise product bound
$Cij$, bandwidth at most $i+j$, and the row/column estimate in
Lemma~\ref{inf:mixed-multipliers}; the available moments absorb
the resulting polynomial factors. The coarser bound $Cp\delta^2$
is therefore valid there as well. Products with one Euler
derivative on either factor obey the same bound in terms of
$\delta$.
All inverse powers of $q,d$ are convergent series in these
multiplier algebras when $C\delta<1/2$.

Subtract $\Delta$ in \eqref{inf:scalar-pullback} and compose with
$\KD$. The four terms have bounds, respectively,
\begin{equation}\label{inf:scalar-four-bounds}
 C\delta,\qquad C\delta,\qquad Cp\delta^2,
                 \qquad C\delta^2(1+\delta).
\end{equation}
For the last bound, \eqref{inf:Cq-identity} gives multipliers
$Cp\delta^2$ and $Cp\delta^3$, and
$\Euler\KD$ costs $C/p$.
The third term uses the bounded operator
$\Euler(\Euler+1)\KD$.
This proves a scalar graph bound $C(\delta+p\delta^2)$ in each
of the three realizations, including the two weighted matrix
norms.

The density requires its own estimate. Both $q$ and $d$ are
harmonic, since $\Delta\Euler q=(\Euler+2)\Delta q=0$.
Write $\alpha=p-1/2$. Direct differentiation gives
\begin{align}
 \frac{\nabla m^{-1}}{m^{-1}}
 &=-\alpha q^{-1}\nabla q-\tfrac12d^{-1}\nabla d,
 \label{inf:density-gradient}\\
 \frac{\Delta m^{-1}}{m^{-1}}
 &=\alpha(\alpha+1)q^{-2}|\nabla q|^2
       +\tfrac34d^{-2}|\nabla d|^2
       +\alpha q^{-1}d^{-1}\nabla q\cdot\nabla d.
 \label{inf:density-Laplacian}
\end{align}
The second formula is quadratic in the shape gradients.
The binomial series for $m^{\pm1}$ has norms at most
$\exp(Cp\delta)$, hence at most two after increasing the
threshold using \eqref{inf:uniform-delta}.
The commutator
\[
 [\Delta,m^{-1}]\KD
 =2\nabla m^{-1}\cdot\nabla\KD+(\Delta m^{-1})\KD
\]
has norm at most
$Cp(\delta/p)+Cp^2(p\delta^2)p^{-2}
=C(\delta+p\delta^2)$ in both weighted matrix algebras.
In particular $H_0m^{-1}H_0^{-1}$ is bounded by an absolute
constant there.

Put $X=(H_\Phi-H_0)H_0^{-1}
=(P_{\Phi_N}\Delta-\Delta)\KD$.
The exact factorization
\[
 \mathcal VH_0^{-1}
 =mX(H_0m^{-1}H_0^{-1})+m[H_0,m^{-1}]H_0^{-1}
\]
therefore has norm at most $C(\delta+p\delta^2)$.
The remaining factor $H_0K_0^{-1}$ is a block-diagonal radial
Hilbert contraction in both matrix norms. This proves the
corresponding estimate for $\mathcal VK_0^{-1}$.
Finally insert \eqref{inf:uniform-delta}, with $\tau\le20$.
The operator constants are absolute; the threshold $P_N$ absorbs
the finite-jet remainder. This proves
\eqref{inf:relative-two-norms} and \eqref{inf:relative-scalar}
with $N$-independent constants, as well as
\eqref{inf:density-graph-bounds}.

We also identify their domains. For this finite smooth chart,
composition and multiplication by the positive smooth density
map the Dirichlet $H^2$ domain to itself. The physical smooth
Dirichlet realization therefore gives the asserted domain and
unitary spectrum. Equivalently the small relative graph bound
in $\Hc$ preserves the ball domain by the closed-operator
perturbation theorem. All equalities above first hold for Poisson
images of polynomial sources. Those images are graph dense in
the Dirichlet domain, because polynomial sources are dense in
$L^2$ and $\KD:L^2\to H^2\cap H^1_0$ is continuous.
Thus the bounded graph extensions are the same geometric
operators. By Lemma~\ref{inf:mixed-domain}, on $\Mmix$ their graph domain is
$\{\KD f:f\in\Mmix\}$ with its source graph norm; the embedding
into $\Hc$ and Poisson uniqueness identify it with the restriction
of that realization. The identical argument applies to $\Cc$.
\end{proof}
